Agents built on large language models (LLMs) now hold conversations that run
for weeks and grow far longer than any context window. On every turn something
has to decide which parts of that history go into the prompt. Re-sending the
whole history makes each reply slower and costlier, a sliding window forgets
the oldest messages whatever they contain, and uniform summarisation loses the
names, numbers and dates that are asked about later. Existing memory systems combine tiers, summaries and forgetting,
but they are evaluated end to end, so the value of each individual decision is
unknown.

This project builds \textbf{MemGate}, a memory layer that sits between an
agent and a frozen LLM and decides, for every message, whether to
\emph{keep} it, \emph{compress} it, \emph{forget} it, or \emph{recall} it
later, without training the base model. It is built as a measurement
harness: a three-tier store runs under a fixed storage budget, and policies
that differ in exactly one factor are compared on the LoCoMo benchmark with
95\% confidence intervals resampled by conversation.

The main findings are as follows. (i) The usual accounting bounds the
context but not the store; once the store is bounded, choosing
\emph{which} turns to forget gains $+9.37$ answer-recall points over
oldest-first eviction at identical storage, fidelity and retrieval, while
compressing evicted turns instead of dropping them loses $7.13$ points.
(ii) No rate--distortion crossover appears up to $182\times$ compression.
(iii) Counting the embedding index in bytes reverses the ranking below about
40\,KiB. (iv) A small LLM judge of salience loses to a hand-written heuristic,
whereas a logistic-regression scorer over shallow features and a sentence
embedding, trained on generated and human conversations, beats the heuristic
by $+14.55$ points on held-out conversations and outperforms rebuilt
MemGPT-style and Mem0-style memory at every budget tested. (v) On held-out
chats of up to ten million tokens it beats retrieval on every corpus, and when
million-token chats are typed into the application it keeps the answer for
51.9\% of questions against 19.2\% for retrieval.

The learned scorer ships in a working chat application that shows, for every
message, the tier it went to and why. The work has been written up as a
research paper with the project mentor as corresponding author.

\vspace{6pt}
\noindent\textbf{Keywords:} large language models, agent memory, context
management, retrieval-augmented generation, learned salience scoring,
rate--distortion, empirical evaluation.
